\section{Indexed data types}
\applabel{indexed data types}

The theory of indexed data types has a few more moving parts. So in
this case we do most of our proofs more abstractly in \appref{cat dev}.
In this appendix we establish the assumptions needed for that development, prove
a result about canonisers, who do not feature in the abstract account, and state
the indexed version of the Parameterised Initiality Theorem.

\begin{proof}[\theoremref{descent}]%
  Let's start with the left adjoint.
  It is an action by direct calculation:
  \[
  \unit \iact\ctx\pair gx \definedby \pair{\unit\mul g}x = \pair gx
  \qquad
  g \act h \act \pair sx \definedby \pair{g\mul(h\mul s)}x
  =
  \pair{(g\mul h)\mul s}x
  =
  (g\mul h)\act \pair sx
  \]
  Given any $\Ctx$-indexed $\G$-action $\A$ and a vertical map $h : \X
  \to \carrier\A$, define:
  \[
  h^{\dagger} : \DPair\g\G{\g\act\X}\equito \A
  \qquad
  h^\dagger_{\ctx}\pair g{x} \definedby \parent{g \iact{\inv g \act\ctx}x} \in \X_{\ctx}
  \]
  Indeed:
  \[
  h^\dagger_{g\act\ctx}(g \act \pair sx) =
  h^\dagger_{g\act\ctx}(\pair{g \mul s}x) =
  (g\mul s)\act (hx)
  =
  g\act (s\act (hx))
  =
  g\act h^\dagger_{\ctx}\pair sx
  \]
  Now, $h^{\dagger}$ extends $h$ along the unit, since:
  \[
  h^{\dagger}\pair\unit x = \unit\act hx = hx
  \]
  If $\overline h : \DPair\g\G{\g\act\X}\equito \A$ is any other such
  extension, we have:
  \[
  \overline h_{\ctx}\pair gx =
  \overline h_{\ctx}(g \act \pair\unit x)
  \explain={equivariance}
  g \act \overline h_{\inv g \act \ctx}\pair\unit x
  \explain*={$\overline h$ extends \\the unit}
  g \act h_{\inv g \act \ctx}x
  =
  h^\dagger_{\ctx}\pair gx
  \]
  and so $\overline h = h^\dagger$, as we wanted.

  The monadicity result follows by an application of Beck's
  monadicity theorem~\cite[Thm.~VI.7.1]{mac-lane:cwm}.
  Start with a split coequalier diagram, in which $\A$ and
  $\B$ are $\Ctx$-indexed actions, and $f,g : \A \equito \B$ are
  equivariant:
  \diagram{28}
  First, equip $\carrier\C$ with an indexed action:
  \(
  h \iact\ctx c \definedby q_{g\act\ctx}(h\iact\ctx s_{\ctx} c)
  \).
  Indeed:
  \(
  \unit\iact\ctx c = q(\unit\iact\ctx s c)
  =
  q(s c) = c
  \)
  and:
  \[
  h \act k \act c = q(h\act s (q(k\act s c)))
  =q(h\act (k\iact\ctx s c))
  =q((h\mul k)\act s c)
  = (h\mul k)\act c
  \]
  The coequaliser $\carrier\Ctx \types q : \carrier\B \to \carrier\C$ in $\Fam_{\carrier\Ctx}$
  is in fact equivariant $\Ctx \types q: \B \equito \C$:
  \begin{align*}
  h\act(q\,b) &= q(h\act(s(q\,b)))= q(h\act(f(t\,b)))
  = q(f(h\act(t\,b)))
  = q(g(h\act (t\,b)))
  = q(h\act g(t\,b))
  \\&= q(h\act b)
  \end{align*}
  It is also a coequaliser in $\Action\G_{\Ctx}$. Given any
  coequalising equivariant map $\ell : \B \equito D$,
  the underlying family map $\ell : \carrier\B \to \carrier D$ is also
  coequalising. Therefore, the unique extension in $\Fam_{\carrier\Ctx}$ of $\ell$
  along $q$ is given by:
  \[
  r : \carrier \C \xto s \carrier\B \xto \ell \carrier D
  \]
  It is an equivariant map $r : \C \equito D$. Indeed:
  \begin{align*}
  r(h \act c) &= \ell(s(h \act c)) = \ell(s(q(h \act s\,c)) =
  \ell(f(t(h \act s\,c)) = \ell(g(t(h\act s\, c)))
  = \ell(h\act s\,c)
  \\&= h \act (\ell(s\,c)) = h \act (r\,c)
  \end{align*}
  Therefore $\carrier-$ lifts split coequalisers of parallel pairs to
  coequalisers.

  Now, if $\C$ already had an indexed $\A$-action structure and $q$
  was equivariant, then necessarily $h \act c = h\act q(s\,c) =
  q(h\act (s\,c))$, and so the action on $\carrier\C$ must be the one
  we defined. Therefore, $\carrier-$ reflects split coequalisers of
  parallel pairs as coequalisers.

  By Beck's theorem, the free indexed action adjunction is monadic.

  Turning to the right adjoint, we first prove it's an action:
  \[
  \unit\iact{\ctx}\phi
  \definedby
  \parent{\fun{h}\phi(h\mul\unit)}
  =
  \parent{\fun{h}\phi h} = \phi
  \] and:
  \begin{align*}
    (g \act k \act \phi)h =
    (k \act \phi)(h \mul g) =
    \phi((h\mul g) \mul k)
    =
    \phi(h\mul (g \mul k))
    =
    ((g\mul k)\act \phi)h
  \end{align*}
  and so \(
  g \act k \act \phi = (g\mul k)\act \phi
  \).

  Next, take any $\Ctx$-indexed $\G$-action $\A$, and a vertical map
  $\carrier\Ctx \types f : \carrier\A \to \X$. Define:
  \[
  \Ctx \types
  f^\dagger : \A \equito \parent{\DFun\g\G{\inv g \act \X}}
  \qquad
  f^\dagger_{\ctx} a \definedby \parent{\fun g f(g \act a)} \in
  \DFun\g{\carrier\G}\parent{\X_{\g \act \ctx}}
  \]
  It is equivariant:
  \[
  f^\dagger_{h\act \ctx}(h \act a) = \fun g f(g\act (h \act a))
  =
  \fun g f((g \mul h)\act a)
  =
  h \act (\fun g f (g \act a))
  =
  h \act f^\dagger_{\ctx} a
  \]

  It lifts $f$ along $\eval\pair{-}\unit$:
  \[
  \eval\pair{f^\dagger a}\unit = (f^\dagger a)\unit = f(\unit \act a) = f\,a
  \]
  Now if $\overline f : \A \equito \DFun\g\G{\inv\g\act \X}$ lifts $f$
  along $\eval\pair{-}\unit$, then:
  \[
  (\overline f a) g = (\overline f a)(\unit \mul g) =
  (g \act (\overline f a))(\unit) =
  \eval\pair{g \act (\overline f a)}\unit
  =
  \eval\pair{\overline f (g \act a)}\unit
  =
  f(g \act a)
  \]
  and so $\overline f = f^\dagger$, and we have a right adjoint.

  We use Beck's theorem in its dual formulation. Consider a split
  equaliser, with $\A$, $\B$ actions:
  \diagram{29}
  Equip $\carrier\C$ with the following indexed $\G$-action structure:
  \(
  k \iact\ctx c \definedby s(k \iact\ctx (q\,c))
  \). It is indeed an action, for $\unit \act c = s(\unit \act (q\,c)) = s(q\,c) = c$
  and:
  \[
  k \act h \act c = s(k \act (q (s(h \act (q\,c))))) =
  s(k \act h \act (q\,c)) = s((k \mul h) \act (q\,c)) = (k \mul h)\act c
  \]
  We need to show the structure maps $q$ and $s$ (but not $t$!)
  are equivariant. For $q$:
  \[
  q(k \act c) = q(s(k \act (q\,c))) = k \act (q\,c)
  \]
  For $s$:
  \[
  k \act s b = s(k \act (q (s\,b))) = s(k\act b)
  \]
  Now, given an equalising equivariant map $\Ctx \types e : \X \equito \B$,
  by moving to families, we know the mediating map is $u : \X
  \xto{e} \B \xto{s} \C$, and therefore equivariant. Therefore
  $\carrier-$ lifts equalisers of split equalisers. Now, equip $\C$ with any
  action that makes $q : \C \to \B$ equivariant. We then have:
  \[
  k\act c = s(q(k\act c)) = s(k\act (q\,c))
  \]
  and so the action on $\C$ is the lifted one, making $q : \C \to \B$
  a coequaliser, i.e., $\carrier-$ reflects equalisers that split in
  $\carrier\Ctx$-families.
\end{proof}

Next, we justify not checking the second axiom of $\G$-strengths in
the cases $g = \unit$ or $h = \unit$:

\begin{proof}[\lemmaref{indexed check only non-trivial}]
  For $h = \unit$:
  \begin{align*}
    \fmap(\mapmul)
  (\gstrength(g,
   \gstrength(\unit, t)))
   &\explain={first axiom}
   \fmap (\mapmul)
  (\gstrength_{\X}(g,
   \fmap(\fun x \pair \unit x) t))
   \\&\explain={naturality}
   \fmap (\mapmul)
   (\fmap(\fun{\pair{g'}x}\triple{g'}\unit x)
     \gstrength(g, t))
   \\&\explain={map fusion}
   (\fmap(\fun{\pair{g'}x}\pair{g'\mul\unit}{x})
   \gstrength(g, t))
   \\&= \gstrength(g, t) = \gstrength(g\mul\unit, t)
  \end{align*}
  Next, for $g = \unit$:
  \begin{align*}
    \fmap(\mapmul)
  (\gstrength(\unit,
   \gstrength(h, t)))
   &\explain={first axiom}
   \fmap(\mapmul)(
   \fmap(\fun{\pair{h'}x}\triple\unit{h'}x)
   (\gstrength(h, t)))
   \\&\explain={map fusion}
   \fmap (\fun{\pair{h'}x}\pair{\unit\mul h'}x)
   (\gstrength(h, t))
   \\&= \gstrength(h, t) = \gstrength(\unit\mul h, t)
  \end{align*}
  as we wanted.
\end{proof}

Next, we state the indexed analogue to the Parameterised Initiality \theoremref{parameterised initiality} for the indexed
setting:
\begin{theorem}[indexed parameterised initiality]%
  \theoremlabel{indexed parameterised initiality}
  Let $\Ctx$ be a $\G$-action, $F : \Fam_{\carrier\Ctx} \to
  \Fam_{\carrier\Ctx}$ an endofunctor equipped with a symmetry
  strength, and $\carrier\Ctx \types \roll : F\fix F \to \fix F$ an
  initial $F$-algebra. For every $F$-algebra $f : F\X \to \X$,
  there is a unique \emph{parameterised homomorphism} into $\pair \X f$, i.e., a
  map
  $\pfold f : \gast \fix F \to \X$ satisfying:
  \(
  (\pfold f)(g, \roll u)
  =
  f\compose \fmap(\pfold f)\strength(g, u)
  \).
\end{theorem}

We will prove this theorem and the Indexed Equivariant Initiality
\theoremref{indexed representation} using the more abstract tools of
the next section.

The remainder of this section establishes the results concerning
canonisers for indexed $\G$-strong functors, which do not have a
general counterpart. By ignoring the indices, they also generalise the
corresponding results from \subsecref{wlog-traversals}.

\begin{lemma}\lemmalabel{canonical elements functor}
  Let $\Ctx$ be a $\G$-action, $F : \Fam_{\carrier\Ctx} \to
  \Fam_{\carrier\Ctx}$ be a functor equipped with a symmetry strength
  $\strength : \gast F \X \to F(\gast \X)$, and $\canoniser : F \to
  \carrier\G$ an indexed canoniser for $F$. The canonical elements extend to a functor
  $\Canon\canoniser : \Action\G_{\carrier\Ctx} \to
  \Fam_{\carrier\Ctx}$ given for $\Ctx \types h : \A \equito \B$ by:
  \[
  ({\Canon\canoniser})\,h\,({g \iact\dtx u}) \definedby
  g \iact\dtx  (F h u)
  \tag{$g = \canoniser_{\carrier\A; \dtx}u, g \act \dtx = \ctx, u \in (F\carrier\A)_{\dtx}$}
  \]
  Consequently, we have an equivariant function $\Ctx \types
  (\parent{\Canon\canoniser}h : (\Canon\canoniser)\A \equito
  (\Canon\canoniser)\A$. The inclusions $\Canon\canoniser \embed F$ are natural.
\end{lemma}
\begin{proof}
  First, since $h : \A \equito \B$ is equivariant, $Fh : F\A \equito F\B$ is
  equivariant. Therefore:
  \[
  (\Canon\canoniser)h(g \iact\dtx u) \definedby
  Fh(g \iact\dtx u) = g \iact\dtx (Fhu)
  \]
  and so the equation in the lemma statement defines
  a family map $k : (\Canon\canoniser)\A \to F\B$,
  independently of the representatives used.

  The crux of the proof is that, by $\canoniser$'s naturality,
  $\canoniser_{\carrier\B; \dtx}(\fmap h u) = \canoniser_{\carrier\A;
    \dtx}(u) = g$. Therefore, $g\iact\dtx(\fmap h u)$ is canonical,
  and $(\Canon\canoniser)h$ is well-defined by lifting $k :
  (\Canon\canoniser)\A \to F\B$ along the inclusion
  $(\Canon\canoniser)\B \embed F\B$. Functoriality follows by routine
  calculation. For equivariance, take:
  \begin{gather*}
    g \iact \dtx u \in ((\Canon\canoniser)\A)_{\ctx}
    \tag{$g = \canoniser_{\carrier\A; \dtx}u, g \act \dtx
      = \ctx, u \in (F\carrier\A)_{\dtx}$}
    \\
    p \iact \xtx v \in ((\Canon\canoniser)\A)_{\ctx}
    \tag{$p = \canoniser_{\carrier\A; \xtx}v, p \act \xtx
      = \ctx, v \in (F\carrier\A)_{\xtx}$}
  \end{gather*}
  and $r \in \stabilise \ctx \subset \carrier\G$ such that:
  \[
  r \iact\ctx g \iact \dtx u = p \iact\xtx v
  \]
  We then have:
  \begin{align*}
  \parent{\Canon\canoniser}h\, (p \iact\xtx v)
  &{}=
  p \iact\xtx (\fmap h\,v)
  \explain={$\fmap h : F\A \equito F\B$}
    \fmap h\,(p \iact\xtx v)
  =
  \fmap h\,(r\iact\ctx g \iact\dtx u)
  =
  r\iact\ctx g \iact\dtx (\fmap h\,u)
  \\
  &{}=
  r\iact\ctx(\parent{\Canon\canoniser}h\, (g \iact\dtx u))
  \end{align*}
  and so $\parent{\Canon\canoniser}h$ is equivariant.

  The naturality of the inclusions $\Canon\canoniser \embed F$ amounts
  to the calculation:
  \[
  (\Canon\canoniser)h(g\iact\dtx u) \definedby g\iact \dtx (Fh u)
  =
  Fh(g \iact \dtx u)
  \]
  which holds by equivariance of $Fh : F\A \equito F\B$.
\end{proof}

The following transformation lets us simultaneously work with
canonical elements in a symmetric fibre while carrying around their
canonising symmetry.

\begin{proposition}
  The transformation $\mctx : \carrier\Ctx \types
  \canonise[\canoniser] \definedby \fun {\BoundV u}
  \pair-{\inv*{\canoniser\BoundV u}}{\canoniser\BoundV u \iact\mctx
    \BoundV u}: F \To \gast (\Canon\canoniser)$ is natural.
\end{proposition}
\begin{proof}
  More generally, for every $g \in \carrier\G$, we have transformation:
  \[
  \ctx : \carrier\Ctx \types (g\iact\ctx\eta-) : \carrier\A \to \gast\carrier\A
  \qquad
  g\iact{\ctx}\eta a \definedby \pair-{g}{\inv g \iact\ctx a}
   = g \iact[\gast\carrier\A]{\inv g\act\ctx} {\pair\unit {\inv g \iact\ctx a}}
  \]
  Given an equivariant family map $\Ctx \types f : \A \equito \B$, we indeed have:
  \diagram{51}
  We now use the canoniser's naturality:
  \begin{align*}
  \canonise(\fmap f u) &{}\definedby
  \pair{\inv*{\canoniser (\fmap f u)}}{(\canoniser (\fmap f u))\act (\fmap f u)}
  \explain={$\canoniser : F \to \carrier\G$ natural}
  \pair{\inv*{\canoniser u}}{(\canoniser u)\act (\fmap f u)}
  \\&\explain={$\fmap f : F \A \equito F\B$}
  \pair{\inv*{\canoniser u}}{\fmap f \parent{(\canoniser u)\act(u)}}
  \explain={$(\inv*{\canoniser u}\act\eta-) : \carrier\C \to \gast\carrier\C$ natural}
  \fmap f \parent{\pair{\inv*{\canoniser u}}{(\canoniser u)\act u}}
  \\&\defines
  \fmap f (\canonise u)
  \end{align*}
  as we wanted.
\end{proof}

We can use this transformation to express the \WLOG{} construction:
\begin{lemma}
  We have:
  \[
  \carrier\Ctx \types
  \wlg(\canoniser, f) : F\carrier\A \xto{\canonise[\canoniser]}
  \carrier{\gast ((\Canon\canoniser)\A)}
  \xto{\gast f} \carrier{\gast\carrier\A}
  \xto{(\act)} \carrier\A
  \]
\end{lemma}
\begin{proof}
  Calculate:
  \begin{align*}
  u \xmapsto{\canonise[\canoniser]}
  \pair{\inv*{\canoniser u}}{(\canoniser u) \act u}
  \xmapsto{\gast f}
  \pair{\inv*{\canoniser u}}{f \parent{(\canoniser u) \act u}}
  \xmapsto{(\act)} {} &
  \inv*{\canoniser u}\act {f \parent{(\canoniser u) \act u}}
  \\&= \wlg(\canoniser, f)
  \end{align*}
  as we wanted.
\end{proof}

The next result establishes a universal property for the intermediate
step in the $\flg$ construction: the extension of a core algebra to an
$F$-algebra.

\begin{theorem}[indexed equivariant algebra representation]%
  \theoremlabel{indexed equivariant algebra representation}
  Let $\Ctx$ be a $\G$-action, $F$ a $\Ctx$-indexed $\G$-strong functor,
  and $\canoniser : F \To
\carrier\G$ an indexed canoniser for $F$, and consider an indexed
equivariant core-$F$ algebra $\Ctx \types f :
(\Canon\canoniser)\A \equito \A$. Then the \WLOG{} construction is the
unique equivariant $F$-algebra $\Ctx \types \wlg(\canoniser, f) :
F\A \equito \A$ extending $f$ along the inclusion
$\Canon\canoniser\embed F$.

Given two equivariant $F$-algebras $(\A, f)$ and $(\B, g)$
and an equivariant function $\Ctx \types h : \A \equito \B$, it is an
equivariant $F$-homomorphism $\Ctx \types h : (\A,f)
\equito (\B, g)$ iff it is an equivariant core-$\Canon\canoniser$
algebra homomorphism.
\end{theorem}
\begin{proof}
  Assume $f : (\Canon\canoniser)\A \equito \A$ is a core algebra.
  By the naturality of $\Canon\canoniser$, $f$ is a core
  function for the canoniser $\canoniser_{\A} : F\carrier A \to
  \carrier\G$. This function is equivariant w.r.t.~$\canoniser_{\A}$
  in the sense of the main paper iff it is equivariant w.r.t.~the
  canoniser $\canoniser$ for $F$ in the sense of \definitionref{indexed core algebras}.

  If $f$ is equivariant, then by the Indexed Representation
  \theoremref{indexed-wlog-construction}, $\wlg(\canoniser, f) = \wlg(\canoniser_{\carrier\A}, f)$ is
  equivariant extending $f$. Covnersely, if $h : F\A
  \equito \A$ is an equivariant algebra, then by the representation
  theorem, $h = \wlg(\canoniser_{\carrier\A}, f) = \wlg(\canoniser,
  f)$, as we wanted.

  Given another such equivariant algebra $g : F\B
  \equito \B$, and an equivariant map $h : \A \to \B$. The naturality
  and injectivity of the embedding $\Canon\canoniser \embed F$ show
  that $h$ is an $F$-algebra homomorphism iff it is a
  $\Canon\canoniser$-algebra homomorphism. Indeed, consider the
  composite square:
  \[
  \diagram*{52}
  \tag{$*$}
  \]
  Our goal is to show that the inner square $(**)$ commutes---$h$ is
  an $F$-algebra homomorphism---iff the outer square ($*$)
  commutes---$h$ is a $\Canon\canoniser$-algebra homomorphism. If
  ($**$) commutes, the whole diagram commutes. Assume ($*$) commutes.
  We have that $f$, $g$, and $h$ are equivariant, and the latter also
  implies that $Fh$ is equivariant. Therefore, the two
  sides of $(**)$ are equivariant extensions of the outer face ($*$)
  along the embedding $(\Canon\canoniser)\A\embed F\A$. By the
  indexed representation theorem, they coincide, i.e., ($**$) commutes.
\end{proof}

We conclude by proving the representation theorem for the indexed
$\flg$ construction:
\begin{proof}
  The proof assembles the previous results. First, by the indexed
  equivariant initiality theorem, the initial algebra is equivariant
  $\roll : F\fix F \equito \fix
  F$. By the algebra representation \theoremref{indexed
    equivariant algebra representation}, its restriction along the
  embedding $\Canon\canoniser \embed F$ is an equivariant algebra
  $\roll : (\Canon\canoniser)\fix F \to \fix
  F$.

  Now, given any equivariant core algebra $f : (\Canon\canoniser) \A
  \equito \A$, by the algebra representation \theoremref{indexed
    equivariant algebra representation}, its \WLOG{} construction is
  an equivariant algebra $g \definedby \wlg(\canoniser, f) :
  F\A \equito \A$.  By the equivariant initiality
  theorem, the unique $F$-homomorphism out of the initial algebra is
  equivariant $\flg f \definedby \fold g : \fix F \equito
  \A$. By the algebra representation \theoremref{indexed equivariant
    algebra representation}, it is also a
  $(\Canon\canoniser)$-homomorphism.

  Conversely, if $h : \fix F \equito \A$ is any
  equivariant core algebra homomorphism, then by \theoremref{indexed
    equivariant algebra representation}, it is an equivariant
  $F$-homomorphism. By the equivariant initiality
  theorem, $h = \fold g = \flg f$.
\end{proof}
